\hypertarget{chapter-03-page-061}{}
\subsection{共轭 Poisson 核}
\label{sec:ch3-conjugate-poisson-kernel}

给定函数 $f\in\mathcal S(\mathbb R)$, 它到上半平面的调和延拓为 $u(x,t)=P_t*f(x)$, 其中 $P_t$ 是 Poisson 核. 由\eqref{eq:ch1-abel-poisson-complex-formula}, 也可写成
\[
 u(z)=\int_0^\infty \widehat f(\xi)e^{2\pi iz\xi}\,d\xi
      +\int_{-\infty}^0 \widehat f(\xi)e^{2\pi i\bar z\xi}\,d\xi,
 \qquad z=x+it.
\]
现在定义
\[
 iv(z)=\int_0^\infty \widehat f(\xi)e^{2\pi iz\xi}\,d\xi
      -\int_{-\infty}^0 \widehat f(\xi)e^{2\pi i\bar z\xi}\,d\xi.
\]
则 $v$ 也在 $\mathbb R_+^2$ 上调和, 而且当 $f$ 为实值函数时, $u$ 和 $v$ 都是实值函数. 此外, $u+iv$ 是解析函数, 因而 $v$ 是 $u$ 的调和共轭.

显然, $v$ 还可写成
\[
 v(z)=\int_{\mathbb R}-i\operatorname{sgn}(\xi)e^{-2\pi t|\xi|}
       \widehat f(\xi)e^{2\pi ix\xi}\,d\xi,
\]
这等价于
\begin{equation}
\label{eq:ch3-conjugate-poisson-convolution}
 v(x,t)=Q_t*f(x),
\end{equation}
其中
\begin{equation}
\label{eq:ch3-conjugate-poisson-fourier}
 \widehat Q_t(\xi)=-i\operatorname{sgn}(\xi)e^{-2\pi t|\xi|}.
\end{equation}
对 Fourier 变换作反演, 得
\begin{equation}
\label{eq:ch3-conjugate-poisson-kernel}
 Q_t(x)=\frac{1}{\pi}\frac{x}{t^2+x^2}.
\end{equation}

\hypertarget{chapter-03-page-062}{}
这就是共轭 Poisson 核. 立即可以验证 $Q(x,t)=Q_t(x)$ 是上半平面内的调和函数, 并且是 Poisson 核 $P_t(x)$ 的调和共轭. 更准确地说,
\[
 P_t(x)+iQ_t(x)=\frac{1}{\pi}\frac{t+ix}{t^2+x^2}=\frac{i}{\pi z},
\]
它在 $\operatorname{Im}z>0$ 时解析.

在第~\ref{chap:hardy-littlewood-maximal-function}~章中, 我们利用 $\{P_t\}$ 是恒等逼近这一事实研究了 $t\to0$ 时 $u(x,t)$ 的极限. 我们也希望对 $v(x,t)$ 作同样处理, 但立即遇到一个障碍: $\{Q_t\}$ 不是恒等逼近, 事实上对任意 $t>0$, $Q_t$ 都不可积. 形式上,
\[
 \lim_{t\to0}Q_t(x)=\frac{1}{\pi x};
\]
它甚至不是局部可积函数, 因而不能定义它与光滑函数的卷积.
